\chapter{Cấu hình tập trung: Config Server}
\label{ch:configserver}

\section{Vấn đề}

Bảy service, mỗi cái đều cần URL database, địa chỉ Kafka, thiết lập JWT,
host mail. Chép những thứ đó vào bảy file \code{application.yml} và mỗi
thay đổi trở thành bảy lần sửa, bảy lần build lại, cùng lời đảm bảo rằng
sớm muộn hai bản sao sẽ lệch nhau. Câu trả lời của Spring Cloud cổ điển:
một HTTP service nhỏ mà công việc duy nhất là trao cho từng ứng dụng cấu
hình của nó lúc khởi động.

\section{Spring Cloud Config hoạt động thế nào}

Server expose cấu hình qua một hợp đồng HTTP đơn giản:

\begin{lstlisting}
GET /{application}/{profile}
GET /course-service/default      -> the config for course-service
\end{lstlisting}

Phía client là dòng \code{spring.config.import:} \code{configserver:...}
bạn đã thấy ở Chương~1. Trong lúc khởi động --- \emph{trước} khi hầu hết
bean được tạo --- client lấy tài liệu của mình về và gộp vào Environment
như một property source đứng trên YAML đóng gói nhưng thua biến môi
trường. Thứ tự phân giải ở Mục~\ref{sec:property-order} làm toàn bộ công
việc; config-server chỉ là thêm một nguồn nữa trong chồng.

\section{Server của dự án, có chú thích}

\begin{lstlisting}[language=yaml]
server:
  port: 8888                       (*@\co{1}@*)
spring:
  application:
    name: config-server
  profiles:
    active: native                 (*@\co{2}@*)
  cloud:
    config:
      server:
        native:
          search-locations: classpath:/configurations   (*@\co{3}@*)
management:
  endpoints:
    web:
      exposure:
        include: health            (*@\co{4}@*)
\end{lstlisting}

\begin{description}
  \item[\co{1}] 8888 là cổng config-server theo quy ước, cũng như 8761 là
  của Eureka. Quy ước tiết kiệm tài liệu.
  \item[\co{2}] Profile \code{native} nghĩa là ``phục vụ file từ một vị
  trí trên filesystem/classpath'' thay vì backend git mặc định.
  \item[\co{3}] Các file được phục vụ nằm \emph{ngay trong jar của chính
  config-server}: \code{configurations/course-service.yml},
  \code{auth-service.yml}, v.v. --- được chọn bằng cách khớp tên file với
  \code{spring.application.name} của service gửi yêu cầu.
  \item[\co{4}] Chỉ \code{health} được expose --- và đây là endpoint mà
  healthcheck của Compose và các probe của Kubernetes dựa vào.
\end{description}

\begin{decision}[classpath native vs.\ backend git]
Backend chuẩn doanh nghiệp là một git repository: thay đổi cấu hình là
các commit --- được kiểm toán, được review, hoàn tác được mà không cần
build lại gì cả. Cái giá là một repository thứ hai và việc quản lý thông
tin đăng nhập. Với một dự án một-lập-trình-viên mà thay đổi cấu hình dù
sao cũng đi cùng thay đổi code, classpath native giữ mọi thứ trong một
repo và một pipeline. Thứ bị đánh đổi: thay một giá trị giờ đòi hỏi build
lại image config-server. Hãy biết nhà tuyển dụng của bạn đứng ở phía nào
của sự đánh đổi đó.
\end{decision}

\begin{pitfall}[service khởi động với giá trị mặc định rồi hỏng một cách kỳ lạ]
Triệu chứng: auth-service khởi động được nhưng mọi request đều thất bại
với lỗi kết nối tới \code{localhost:5432}. Nguyên nhân: config-server
không liên lạc được lúc khởi động, và \code{optional:configserver:} để
service khởi động từ giá trị mặc định đóng gói sẵn --- vốn trỏ về
localhost. Vì thất bại xuất hiện vài phút sau và cách xa nguyên nhân, nó
trông như một vấn đề database. Kiểm tra: log khởi động in ra cấu hình từ
xa có được lấy về hay không. File Compose ngăn vấn đề thứ tự này bằng
\code{depends\_on: condition: service\_healthy}; Kubernetes không có thứ
tự khởi động (Chương~15), nên các service phải chịu được một
config-server đến muộn bằng cách crash rồi khởi động lại cho đến khi nó
xuất hiện.
\end{pitfall}

\section{Thứ tự khởi động và điệu nhảy bootstrap}

Config-server tạo ra một ràng buộc trình tự ngầm trên toàn hệ thống: mọi
service khác đều muốn nó lên trước. Bạn sẽ gặp hệ quả của điều này ba lần
trong cuốn sách:

\begin{itemize}
  \item Compose mã hóa thứ tự một cách tường minh bằng \code{depends\_on}
  và một healthcheck dựa trên \code{curl} (Chương~12).
  \item Kubernetes từ chối mã hóa thứ tự khởi động; hệ thống hội tụ bằng
  retry (Chương~15).
  \item Pipeline Jenkins triển khai mọi manifest trong một lệnh
  \code{kubectl apply -k} và để các rollout tự ổn định đồng thời
  (Chương~23).
\end{itemize}

\section{Làm mới cấu hình mà không khởi động lại}

Mặc định, một client đọc cấu hình của nó một lần lúc khởi động. Spring
Cloud cung cấp \code{/actuator/refresh} (lấy lại theo yêu cầu, từng
instance, áp dụng cho các bean \code{@RefreshScope}) và Spring Cloud Bus
(phát tán lệnh refresh qua một message broker). Dự án này không dùng cả
hai: trên Kubernetes một thay đổi cấu hình được giao đi dưới dạng image
mới hoặc thay đổi env, và rolling update của Deployment (Chương~15) dù
sao cũng khởi động lại các pod. Đó là câu trả lời hiện đại cho câu hỏi mà
config-server ban đầu được thiết kế để giải quyết.

\section{Bên ngoài dự án}

Một thiết lập doanh nghiệp điển hình dựa trên git, để đối chiếu:

\begin{lstlisting}[language=yaml]
spring:
  cloud:
    config:
      server:
        git:
          uri: https://git.company.com/platform/config-repo
          default-label: main
          search-paths: '{application}'
\end{lstlisting}

Cấu hình nằm theo từng service trong các thư mục của một repo riêng;
label là một branch, nên \code{default-label: release-42} ghim một môi
trường vào một bản phát hành. Mẫu phổ biến thứ hai bỏ hẳn config-server:
ConfigMap và Secret của Kubernetes được gắn vào dưới dạng biến môi trường
--- chính là hướng dự án này đang đi tới (Chương~16 và 26); config-server
tồn tại ở đây chủ yếu như một ví dụ sống của stack cổ điển.

\begin{quiz}
\begin{enumerate}
  \item Các nguồn này thắng theo thứ tự nào cho
  \code{spring.datasource.url}: YAML đóng gói, tài liệu config-server,
  biến môi trường của Deployment?
  \item \code{profiles.active: native} thay đổi gì về nơi server tìm tài
  liệu, và backend git mua được gì mà native không thể?
  \item Vì sao \code{optional:} trong \code{spring.config.import} là con
  dao hai lưỡi? Mô tả kiểu thất bại mà nó cho phép.
  \item Cấu hình của chính config-server không thể đến từ một
  config-server. Nó phải nằm ở đâu, và đó là nguyên lý bootstrap tổng
  quát nào?
\end{enumerate}
\end{quiz}
